استمر مؤشر أداء الأعمال في الارتفاع خلال الربع محل الدراسة (أكتوبر-ديسمبر 2025) ليتجاوز \textbf{المستوى المحايد} بـ 10نقاط محققا أداء أفضل من الربعين السابق والمناظر. ويرجع ذلك إلى تعافي مؤشرات الإنتاج، والمبيعات المحلية، والصادرات، واستغلال الطاقة الإنتاجية، والتي سجلت جميعها قيما أعلى من \textbf{المستوى المحايد}، وأعلى من مثيلها في الربعين السابق والمناظر، وعلى مستوى كافة الشركات. ويعكس هذا التعافي عدة عوامل أبرزها الاستقرار النسبي للأوضاع الجيوسياسية في الشرق الأوسط، إضافة إلى انتعاش الأسواق الخارجية لصالح المنتجات المصرية، وخاصة بعد فرض الرسوم الجمركية الجديدة، حيث شهدت الفترة محل الدراسة موسمية سواء لموسم رمضان والأعياد (الشكل 1-1).
